% ============================================================
%  Schemes, not Varieties
%  TU Berlin, Summer Semester 2025
% ============================================================

\documentclass[a4paper, 12pt]{article}

% ---- Font & layout (document-specific) ----------------------
\usepackage{mathpazo}
\linespread{1.05}
\usepackage{geometry}
\geometry{margin=2.5cm, top=3cm, bottom=3cm}
\setlength{\parskip}{0.4em}
\setlength{\parindent}{1.5em}

% ---- Section headings ---------------------------------------
\usepackage{titlesec}
\titleformat{\section}
  {\large\bfseries}{\thesection.}{0.6em}{}
  [\vspace{0.3ex}\rule{\linewidth}{0.4pt}]
\titlespacing*{\section}{0pt}{2.5ex}{1.2ex}

% ---- Running headers ----------------------------------------
\usepackage{fancyhdr}
\pagestyle{fancy}
\fancyhf{}
\lhead{\small\itshape Schemes, not Varieties}
\rhead{\small\thepage}
\renewcommand{\headrulewidth}{0.4pt}
\fancypagestyle{plain}{%
  \fancyhf{}
  \rhead{\small\thepage}
  \renewcommand{\headrulewidth}{0pt}
}

% ---- Shared preamble ----------------------------------------
% ============================================================
%  Shared preamble — \input this from every notes document.
%  Do NOT put \documentclass, fonts, geometry, titlesec,
%  or fancyhdr here; those stay per-document.
% ============================================================

% ---- Encoding & typography ----------------------------------
\usepackage[utf8]{inputenc}
\usepackage[T1]{fontenc}
\usepackage{microtype}

% ---- Mathematics --------------------------------------------
\usepackage{amsmath, amsthm, amssymb}
\usepackage{mathtools}

% ---- Lists --------------------------------------------------
\usepackage{enumitem}


% ---- Colour & links -----------------------------------------
\usepackage{xcolor}
\usepackage{hyperref}
\hypersetup{colorlinks=true, linkcolor=black!60, citecolor=black!60, urlcolor=black!60}

% ---- Diagrams -----------------------------------------------
\usepackage{tikz-cd}

% ---- Theorem environments -----------------------------------
\theoremstyle{definition}
\newtheorem{example}{Example}
\newtheorem{remark}[example]{Remark}
\newtheorem{definition}[example]{Definition}

\theoremstyle{plain}
\newtheorem{theorem}[example]{Theorem}
\newtheorem{proposition}[example]{Proposition}
\newtheorem{lemma}[example]{Lemma}
\newtheorem{corollary}[example]{Corollary}
\newtheorem{claim}[example]{Claim}

% ---- Operators ----------------------------------------------
\DeclareMathOperator{\Spec}{Spec}
\DeclareMathOperator{\Proj}{Proj}
\DeclareMathOperator{\Gal}{Gal}
\DeclareMathOperator{\GL}{GL}
\DeclareMathOperator{\Lie}{Lie}
\DeclareMathOperator{\Tor}{Tor}
\DeclareMathOperator{\Char}{char}
\DeclareMathOperator{\height}{ht}
\DeclareMathOperator{\V}{\textbf{V}}

\newcommand{\Aut}{\mathrm{Aut}}
\newcommand{\Hom}{\mathrm{Hom}}
\newcommand{\Ext}{\mathrm{Ext}}
\newcommand{\EExt}{\mathcal{E}xt}
\newcommand{\Def}{\mathrm{Def}}
\newcommand{\Supp}{\mathrm{Supp}}

% ---- Blackboard bold ----------------------------------------
\renewcommand{\AA}{\mathbb{A}}
\newcommand{\GG}{\mathbb{G}}
\newcommand{\PP}{\mathbb{P}}
\newcommand{\QQ}{\mathbb{Q}}
\newcommand{\RR}{\mathbb{R}}
\newcommand{\CC}{\mathbb{C}}
\newcommand{\NN}{\mathbb{N}}
\newcommand{\ZZ}{\mathbb{Z}}

% ---- Calligraphic / fraktur ---------------------------------
\newcommand{\cO}{\mathcal{O}}   % structure sheaf
\newcommand{\cM}{\mathcal{M}}
\newcommand{\cB}{\mathcal{B}}
\newcommand{\cT}{\mathcal{T}}
\newcommand{\fX}{\mathfrak{X}}

% ---- Misc ---------------------------------------------------
\newcommand{\kk}{k}


% ---- Title --------------------------------------------------
\title{\textbf{Schemes, not Varieties}}
\author{Dominic Bunnett}
\date{Last updated: \today}

\begin{document}
\maketitle
\tableofcontents
\newpage

% =============================================================
\section*{Introduction}

These are collective notes on phenomena in scheme theory with no analogue
in classical algebraic geometry over an algebraically closed field.

% =============================================================

\begin{example}[A morphism of schemes, not of varieties]
Let $k$ be a field with algebraic closure $\bar{k}$, and let
$\sigma \in \Gal(\bar{k}/k)$ be non-trivial. Consider $X = \Spec(\bar{k})$.
\begin{enumerate}
\item \emph{Topologically trivial.} $\Spec(\bar{k})$ is a single point, so any
  self-map is the identity on the underlying topological space.
\item \emph{Non-trivial as a scheme morphism.} $\sigma$ induces a ring automorphism
  $\sigma^* \colon \bar{k} \xrightarrow{\sim} \bar{k}$, giving a non-trivial
  automorphism $(\mathrm{id}, \sigma^*)$ of $X$ as a $k$-scheme. In particular
  $\mathrm{Aut}_{k\text{-}\mathrm{sch}}(\Spec \bar{k}) \cong \Gal(\bar{k}/k)$.
\item \emph{Concretely.} Complex conjugation is an order-$2$ automorphism of
  $\Spec(\CC)$ as an $\RR$-scheme: invisible on points, non-trivial on the
  structure sheaf via $a + bi \mapsto a - bi$.
\item \emph{Moral.} A morphism of schemes is a pair (continuous map, morphism of
  sheaves). Classical geometry sees only the first component.
\end{enumerate}
\end{example}

% =============================================================

\begin{example}[{$R = k[x,y]/(xy)$, two lines meeting at a point}]
The coordinate ring of the union of the coordinate axes in $\AA^2$.
This is still a variety — just a reducible one, the union of two lines.
Nothing genuinely new happens here from the scheme-theoretic point of view;
it is included as a baseline before the more exotic examples.
\begin{enumerate}
\item \emph{Reduced, not a domain.} $\bar x \cdot \bar y = 0$ with
  $\bar x, \bar y \neq 0$, but no nilpotents.
\item \emph{Minimal primes.} $(\bar x)$ and $(\bar y)$, one for each irreducible
  component. $\Spec R$ is two copies of $\AA^1_k$ glued at the origin —
  entirely visible to classical geometry.
\item \emph{Not regular at the origin.} The local ring at
  $\mathfrak{m} = (\bar x, \bar y)$ is not a DVR:
  $\mathfrak{m}/\mathfrak{m}^2$ has dimension $2$ over $k$
  while $\dim_{\mathfrak{m}} R = 1$.
\end{enumerate}
\end{example}

% =============================================================

\begin{example}[{$R = k[x]_{(x)}$, localization at a point}]
The localization of $k[x]$ at the prime ideal $(x)$.
\begin{enumerate}
\item \emph{Local.} The unique maximal ideal is $(x)R$, with residue field $k$.
\item \emph{Noetherian.} Localization of a Noetherian ring.
\item \emph{Dimension $1$.} The chain $(0) \subset (x)R$ is maximal.
\item \emph{Regular.} $R$ is a DVR with uniformiser $x$.
\item \emph{$\Spec R$.} Two points: the closed point $(x)$ and the generic
  point $(0)$. Geometrically this is the germ of $\AA^1_k$ at the origin —
  a scheme with no analogue in classical geometry, which sees only closed points.
\end{enumerate}
\end{example}

% =============================================================

\begin{example}[{$R = k[x]_{(x)} \cap k[x]_{(x-1)} \subseteq k(x)$, semilocalization at two points}]
The ring of rational functions regular at both $0$ and $1$.
\begin{enumerate}
\item \emph{Not local.} Two maximal ideals, corresponding to $x = 0$ and $x = 1$.
\item \emph{Noetherian.} A localization of $k[x]$.
\item \emph{Dimension $1$.}
\item \emph{$\Spec R$.} Three points: the generic point $(0)$ and two closed
  points. This is the open subscheme of $\AA^1_k$ obtained by discarding all
  closed points except $\{0, 1\}$ — a purely scheme-theoretic construction.
\end{enumerate}
\end{example}

% =============================================================

\begin{example}[$R = k(x)$, the generic point]
The rational function field in one variable.
\begin{enumerate}
\item \emph{$\Spec R$ is a point.} $k(x)$ is a field, so $(0)$ is its only prime.
\item \emph{Local and regular.} Fields are regular local rings of dimension $0$.
\item \emph{The morphism $\Spec k(x) \to \Spec k[x]$.} Maps the unique point to
  the generic point $(0) \in \Spec k[x]$. The image is the non-closed point,
  invisible in classical geometry over $k$.
\end{enumerate}
\end{example}

% =============================================================

\begin{example}[{$R = k[x_1, x_2, x_3, \ldots]$, infinite affine space}]
The polynomial ring in countably many variables.
\begin{enumerate}
\item \emph{Not Noetherian.} The chain $(x_1) \subset (x_1, x_2) \subset \cdots$
  does not stabilize.
\item \emph{Infinite Krull dimension.} The chain
  $(0) \subset (x_1) \subset (x_1,x_2) \subset \cdots$
  is a strictly increasing chain of primes, so $\dim R = \infty$.
\item \emph{Reduced domain.} $R$ is an integral domain with no nilpotents.
\end{enumerate}
\end{example}

% =============================================================

\begin{example}[{$R = k[[x]]$, formal power series}]
\begin{enumerate}
\item \emph{Local.} Unique maximal ideal $(x)$, residue field $k$.
\item \emph{Noetherian.} Every ideal is of the form $(x^n)$.
\item \emph{Dimension $1$ and regular.} $R$ is a DVR, like $k[x]_{(x)}$.
\item \emph{$\Spec k[[x]]$ vs $\Spec k[x]_{(x)}$.} Both have the same underlying
  topological space (two points), but the local rings differ — one is algebraic,
  the other formal. This shows that two schemes can be homeomorphic but not isomorphic.
\end{enumerate}
\end{example}

% =============================================================

\begin{example}[{$R = k[[x,\varepsilon]]/(\varepsilon^2)$, first-order thickening}]
\begin{enumerate}
\item \emph{Local.} Maximal ideal $(\bar x, \bar\varepsilon)$.
\item \emph{Not reduced.} Nilradical $(\bar\varepsilon)$, since $\varepsilon^2 = 0$.
\item \emph{Dimension $1$.} Nilpotents do not contribute to chains of primes.
\item \emph{Geometric meaning of $\varepsilon$.} $\Spec R$ is a first-order thickening
  of $\Spec k[[x]]$: the extra data of $\varepsilon$ records a tangent direction.
  A map $\Spec R \to X$ is an arc on $X$ together with a first-order deformation —
  the basic object of deformation theory.
\end{enumerate}
\end{example}

% =============================================================

\begin{example}[{$R = \ZZ[x]$, the arithmetic line}]
\begin{enumerate}
\item \emph{Noetherian, dimension $2$.} The chain
  $(0) \subset (p) \subset (p, x-a)$ is maximal, giving $\dim \ZZ[x] = 2$.
\item \emph{Fibres of $\Spec \ZZ[x] \to \Spec \ZZ$.} Over $(p)$:
  $\Spec \mathbb{F}_p[x] = \AA^1_{\mathbb{F}_p}$. Over $(0)$:
  $\Spec \QQ[x] = \AA^1_\QQ$. The arithmetic line interpolates between
  all characteristics simultaneously.
\item \emph{Jacobson.} Every prime is an intersection of maximal ideals.
  Compare $k[x]_{(x)}$, which is not Jacobson: localization kills closed points
  and destroys the Jacobson property.
\end{enumerate}
\end{example}

% =============================================================

\begin{example}[$A^+$, absolute integral closure]
Let $A$ be a domain with fraction field $K$; define $A^+$ as the integral closure
of $A$ in an algebraic closure $\overline K$.
\begin{enumerate}
\item \emph{Domain.} Integral closure of a domain is a domain.
\item \emph{Not Noetherian.} For $A = \ZZ$, $A^+ = \overline{\ZZ}$ (ring of all
  algebraic integers) has infinite Krull dimension.
\item \emph{Not finite over $A$.} Requires infinitely many generators over $A$
  in general.
\item \emph{Relevance.} Appears naturally in perfectoid geometry and $p$-adic
  Hodge theory, where $\ZZ^+_p$ plays a central role.
\end{enumerate}
\end{example}

% =============================================================

\begin{example}[{$R = k[x^{1/p^\infty}]$, perfection in characteristic $p$}]
Let $k$ have characteristic $p > 0$; set
$R = k[x, x^{1/p}, x^{1/p^2}, \ldots]$.
\begin{enumerate}
\item \emph{Domain, not Noetherian.} The chain
  $(x) \subset (x^{1/p}) \subset (x^{1/p^2}) \subset \cdots$
  does not stabilize.
\item \emph{Frobenius is an isomorphism.} $x^{1/p^n} \mapsto x^{1/p^{n-1}}$
  shows $F$ is surjective; it is injective since $R$ is a domain.
  This is what it means for $R$ to be \emph{perfect}.
\item \emph{Dimension $1$.} Same prime structure as $k[x]$ but with a
  richer ring of global functions.
\end{enumerate}
\end{example}

% =============================================================

% =============================================================
\bibliographystyle{alpha}
\begin{thebibliography}{9}

\bibitem{Mumford}
  D.~Mumford,
  \emph{The Red Book of Varieties and Schemes},
  Lecture Notes in Mathematics 1358,
  Springer, 2nd ed., 1999.

\end{thebibliography}

\end{document}
